\chapter{Problèmes de synthèse}
Les synthèses S1 à S4 et S8 mobilisent la spécialité et ses approfondissements signalés ; S5 à S7 mobilisent aussi les mathématiques expertes.
\section{S1 -- Héron : convergence et efficacité}
\textbf{Stratégie :} commencer par un invariant de positivité, puis utiliser l'ordre avant de résoudre l'équation de limite.
Tous les termes sont positifs par récurrence. Pour $u>0$, $\frac12(u+a/u)-\sqrt a=(u-\sqrt a)^2/(2u)\ge0$ ; chaque terme après le premier majore donc la racine.
Dès $n\ge1$, $u_{n+1}-u_n=(a-u_n^2)/(2u_n)\le0$. La suite décroît et est minorée par $\sqrt a>0$, donc converge vers $L>0$. La continuité sur $]0,+\infty[$ permet $L=(L+a/L)/2$, d'où $L^2=a$ et $L=\sqrt a$.

L'identité du premier paragraphe donne aussi l'erreur exacte demandée. Pour $a=2$, les termes sont $u_1=3/2$, $u_2=17/12$, $u_3=577/408$. Comme $u_n\ge\sqrt a$, on a $a/u_n\le\sqrt a\le u_n$. Ainsi
\[
\frac{816}{577}\le\sqrt2\le\frac{577}{408},\qquad
\frac{577}{408}-\frac{816}{577}=\frac1{235416}.
\]
Le milieu approche la racine avec erreur au plus $1/470832$. Arrêter lorsque $u_n-a/u_n\le2\varepsilon$ certifie une erreur de milieu au plus $\varepsilon$, sous calcul exact ou contrôle d'arrondi.

Newton pour $f(x)=x^2-a$ donne $x-f(x)/f'(x)=(x+a/x)/2$. La positivité et la forme particulière ont permis la preuve ; d'autres fonctions peuvent produire des cycles, comme au chapitre 24. L'erreur au carré explique l'efficacité locale de Héron, mais les flottants finissent par limiter la précision.

\section{S2 -- Géométrie d’un tétraèdre}
Les intercepts donnent $x/2+y/3+z/6=1$, soit $3x+2y+z=6$. Les trois points satisfont cette équation et ne sont pas alignés ; elle définit donc leur plan. Un normal est $n=(3,2,1)$.
La droite normale issue de $O$ est $(x,y,z)=(3t,2t,t)$. Dans le plan, $14t=6$, donc $H=(9/7,6/7,3/7)$ et $OH=6/\sqrt{14}$ cm. Le signe positif du paramètre correspond au plan situé dans la direction du normal choisi.

Le triangle rectangle $OAB$ a aire $2\times3/2=3\,\mathrm{cm}^2$. La hauteur issue de $C$ vaut $6$ cm ; le volume vaut $3\times6/3=6\,\mathrm{cm}^3$. En prenant $ABC$ pour base, $6=\operatorname{aire}(ABC)\,OH/3$, donc l'aire vaut $3\sqrt{14}\,\mathrm{cm}^2$.

La sphère a centre $(0,0,3)$ et rayon $3$ ; développer $x^2+y^2+(z-3)^2=9$ donne l'équation. Dans le plan $z=3$, elle devient $x^2+y^2=9$ : un cercle de centre $(0,0,3)$ et rayon $3$. Les équations établissent les appartenances et les distances ; une perspective peut les suggérer mais déforme angles et longueurs.

\section{S3 -- Logarithme, intégrale et sommes}
La fonction $x\mapsto1/(1+x)$ décroît sur $[0,1]$. Les rectangles à droite minorent son intégrale, ceux à gauche la majorent. Son intégrale vaut $[\ln(1+x)]_0^1=\ln2$. En soustrayant les sommes, les termes intérieurs disparaissent et $L_n-R_n=(1-1/2)/n=1/(2n)$.
Le milieu a donc erreur au plus $1/(4n)$. Choisir $n=2500$ suffit pour $10^{-4}$ ; le programme calcule les deux sommes ou l'une puis ajoute leur écart. L'arrondi flottant doit être distingué de cette erreur théorique.

Comme $0\le1/(1+x)\le1$, $0\le J_n\le\int_0^1x^n=1/(n+1)$. L'identité $x^n/(1+x)+x^{n+1}/(1+x)=x^n$ donne la récurrence. En partant de $J_0=\ln2$, elle conduit par récurrence à
\[
\ln2=\sum_{k=1}^{n}\frac{(-1)^{k-1}}k+(-1)^nJ_n.
\]
Le reste a module au plus $1/(n+1)$, et son signe est connu. La limite de la somme alternée vaut donc $\ln2$. Les rectangles fournissent ici une meilleure borne à nombre de termes comparable, mais d'autres algorithmes, dont Briggs, peuvent converger plus vite. Une impression décimale n'est un certificat que si l'erreur totale est contrôlée.

\section{S4 -- Surbooking : modèle et décision}
$X_n$ est somme de $n$ Bernoulli indépendantes : $X_n\sim\mathcal B(n,0{,}9)$, $E(X_n)=0{,}9n$, $V(X_n)=0{,}09n$. Le risque exact est
$\sum_{k=101}^n\binom nk0{,}9^k0{,}1^{n-k}$, nul si $n=100$.
Partir de $n=100$, calculer le risque pour $n+1$, avancer tant qu'il est au plus $0{,}05$, puis retourner le dernier entier admissible. Un calcul rationnel ou des probabilités numériques stables permettent de vérifier la frontière. Le programme livré trouve $n=106$ ; les risques pour $106$ et $107$ valent environ $0{,}0398$ et $0{,}0806$.

Coupler les variables par $X_{n+1}=X_n+B_{n+1}$ avec $B_{n+1}\ge0$ montre l'inclusion $\{X_n>100\}\subset\{X_{n+1}>100\}$. Pour $0{,}9n>100$, Tchebychev donne $P(X_n\le100)\le0{,}09n/(0{,}9n-100)^2\to0$ ; la recherche finit donc par dépasser le seuil.
Pour $n=105$, $E=94{,}5$, $V=9{,}45$. L'événement $X_n>100$ implique $X_n\ge101$, donc un écart au moins $6{,}5$. La borne est $9{,}45/6{,}5^2\simeq0{,}224$, trop grande pour décider la contrainte. Une borne supérieure au seuil ne prouve pas que le risque le dépasse.

Le modèle suppose une probabilité identique et une indépendance entre personnes. Groupes, météo, annulations communes et estimation de $0{,}9$ peuvent invalider ces hypothèses. Le calcul n'inclut ni les conséquences du refus d'entrée ni les contraintes pratiques : il illustre un raisonnement, sans constituer une politique réelle.

\section{S5 -- Complexes, rotations et polygones}
Si $\omega^j=\omega^k$, alors $(j-k)/5$ est entier ; avec $0\le j,k<5$, cela impose $j=k$. Tous les points sont sur le cercle unité et les angles successifs valent $2\pi/5$ : le polygone est régulier. La somme géométrique est $(1-\omega^5)/(1-\omega)=0$, donc le barycentre est l'origine.

$1-e^{it}=e^{it/2}(e^{-it/2}-e^{it/2})=-2ie^{it/2}\sin(t/2)$ ; prendre les modules donne la formule. Le côté vaut $2\sin(\pi/5)$, la diagonale $2\sin(2\pi/5)$.
Diviser $1+\omega+\omega^2+\omega^3+\omega^4=0$ par $\omega^2$ donne $1+s+(s^2-2)=0$. Or $s=2\cos(2\pi/5)>0$, car $0<2\pi/5<\pi/2$. Donc $s=(-1+\sqrt5)/2$ et $\cos(2\pi/5)=(\sqrt5-1)/4$.

La matrice est $R=\begin{pmatrix}\cos(2\pi/5)&-\sin(2\pi/5)\\\sin(2\pi/5)&\cos(2\pi/5)\end{pmatrix}$. Composer cinq rotations additionne les angles ; $R^5=R_{2\pi}=I$. C'est aussi la traduction de $\omega^5=1$ par l'identification du plan à $\mathbb C$.

\section{S6 -- RSA : preuve et limites}
Euclide donne $40=13\times3+1$, donc $1=40-13\times3$ et $d\equiv-13\equiv27\pmod{40}$. Ainsi $ed=81$.
Le chiffré vaut $7^3=343\equiv13\pmod{55}$. Pour déchiffrer, les carrés successifs donnent $13^2\equiv4$, $13^4\equiv16$, $13^8\equiv36$, $13^{16}\equiv31$. Comme $27=16+8+2+1$, le produit $31\times36\times4\times13$ est congru à $7$ modulo $55$.

Si $5\mid m$, $m^{81}\equiv m\equiv0\pmod5$. Sinon Fermat donne $m^4\equiv1\pmod5$, donc $m^{81}=m(m^4)^{20}\equiv m$. De même modulo $11$ : si $11\mid m$, résultat immédiat ; sinon $m^{10}\equiv1$ et $81=1+8\times10$. Ainsi $5$ et $11$ divisent $m^{81}-m$. Étant premiers entre eux, leur produit divise cette différence par Gauss.

L'exponentiation rapide évite les puissances immenses ; la preuve de correction ne prouve pas la sécurité. Les facteurs minuscules sont immédiatement retrouvés, et le RSA brut sans encodage approprié possède d'autres faiblesses. L'exemple sert à relier congruences, Bézout et Fermat ; les protocoles cryptographiques réels demandent une étude spécialisée.

\section{S7 -- Trois états et convergence}
Les coefficients sont positifs et chaque ligne somme à $1$. Les colonnes somment aussi à $1$, donc $\pi P=\pi$. Chaque coefficient de $J^2$ vaut $3$, donc $J^2=3J$ et $Q^2=Q$. La décomposition annoncée donne des coefficients diagonaux $1/3+(1/4)(2/3)=1/2$ et hors diagonale $1/3-(1/4)(1/3)=1/4$.

On a $Q(I-Q)=(I-Q)Q=0$ et $(I-Q)^2=I-Q$. Par récurrence, les produits mixtes s'annulent et $P^n=Q+4^{-n}(I-Q)$, y compris $n=0$ où la somme vaut $I$.
Toute distribution ligne $\mu_0$ satisfait $\mu_0Q=\pi$ ; ainsi
$\mu_n=\pi+4^{-n}(\mu_0-\pi)\to\pi$.
Si $\nu$ est invariante, $\nu=\nu P^n$ pour tout $n$, donc le passage à la limite donne $\nu=\pi$ : unicité démontrée.

Chaque coordonnée de $\mu_0$ appartient à $[0,1]$, donc son écart à $1/3$ est au plus $2/3$. Il suffit de $(2/3)4^{-n}\le10^{-3}$, soit $4^n\ge2000/3$. Comme $4^4=256$ et $4^5=1024$, $n=5$ convient. Ce contrôle quantifie la convergence ; il ne se limite pas à l'observer sur quelques itérations.

\section{S8 -- Une fonction définie par une intégrale}
L'intégrande est continue, donc le théorème fondamental donne $F'(x)=e^{-x^2}>0$ ; la fonction croît strictement. De plus $F''(x)=-2xe^{-x^2}\le0$ pour $x\ge0$, donc elle est concave.
La tangente à l'exponentielle donne $e^{-u}\ge1-u$. Pour $0\le t\le x\le1$, $1-t^2\le e^{-t^2}\le1$. Intégrer donne $x-x^3/3\le F(x)\le x$.

Pour $x\ge1$, $F(x)=F(1)+\int_1^xe^{-t^2}\,dt\le1+\int_1^xe^{-t}\,dt\le1+e^{-1}$. Une fonction croissante majorée possède une limite finie (prendre le supremum de ses valeurs et utiliser la même preuve que pour les suites). Cette limite est positive, puisque $F(1)\ge2/3$.
Pour $y\ge x\ge1$, $0\le F(y)-F(x)\le\int_x^ye^{-t}\,dt\le e^{-x}$. Faire tendre $y$ vers l'infini donne le contrôle du reste.

Une intégrale définit une fonction même lorsqu'aucune expression avec les fonctions élémentaires usuelles n'est disponible. Nous n'avons pas démontré ici une impossibilité d'expression, ni évalué $L$ : sa valeur $\sqrt\pi/2$ nécessite une méthode supplémentaire, étudiée en licence. Après normalisation, $e^{-t^2}$ intervient dans une densité gaussienne ; il faudra définir probabilités continues et intégrales impropres pour rendre ce lien complet.
